% !TeX root = ../papers/nc-p3.tex
\section{高階数成分の全二十印の分離、周期と有理完全不変量}
\label{sec:marked-twenty-rational-invariants}
\begin{theorem}\label{thm:marked-twenty-rational-invariants}
分裂混合点を通る十五次元成分の非空な稠密開部分では、全二十印で
交叉交換子の両側の線型復元が成立する。
基点を固定した添字保持同型の完全不変量は有理演算だけで求まり、
無印の添字保持同型も一つのモニック二十次式と次数20未満の係数多項式の一致で判定できる。
同じ開部分で射影自己同型群は自明、最小周期は4である。
これらの結論には以下の付値環で全次数正則な実現がある。
素数1009で指定した三十四自由係数の合同近傍の\emph{全て}の持上げが
標準Koszul左右AS正則であり、そこに互いに添字保持非同型な十五パラメータの解析的族を取れる。
代数的数係数でのそのような正則閉点も稠密に存在する。
\end{theorem}

\begin{proof}[二個の行列が与える有理な基底]
前節の $C:\mathfrak{sl}(V)\to\mathfrak{sl}(U)$ を使う。
$T=\operatorname{Tr}(C^*C)\ne0$ とし、両側で $K=C^*C/T,CC^*/T$ を取る。
$\operatorname{Tr}$ は八次元作用の跡であり、以下の $\operatorname{tr}$ は三次行列の跡である。
跡対合の互いに双対な基底 $e_i,e^i$ について、三次行列
\[
 J(H)=\Big(\sum_i H(e_i)e^i+e^i H(e_i)\Big)_0,\quad
 A=J(K),\quad B=J(K^2),\quad m=\operatorname{tr}(A^2),\quad d=\det[A,B]
\]
を作る。添字0は跡零部である。双対基底の縮約だから $J$ は基底の選択に依存せず、
各側の共役と可換である。$K$ は $C$ の全体倍率にも不変である。
入力側を $V$、出力側を $U$ と添字付けし、
$\tau=\operatorname{tr}(A_U C(A_V))$ が非零なら実際の $C$ を $C/\tau$ で規格化する。
次の九語が各側で全行列環の基底である開条件を課す。
\[
           I,A,B,A^2,AB,BA,B^2,A^2B,AB^2 .
\]
各基底での行列積表と、単位を零に送って延長した $C/\tau$ の九語間の線型写像の表を $H_p$
とする。どの表も固定順の基底であり、有理演算だけで得られる。
等しい積表から全行列環の単位的同型が二つ得られ、いずれも内部共役である。
第三の表との一致で $[C]$ を運ぶので、前節の両側単射条件の下でこの表は完全である。
また自己同型では各側の $A,B$ が個別に固定され、九語が基底なので両共役とも恒等である。
残るスカラーは前節の逆倍率の項で尽き、印を固定する射影自己同型はない。
\end{proof}

\begin{proof}[二十個の印の区別と四つの開始位置]
印 $p$ を一つ先に運んだ位置で、元の星係数から読む四組は
$(G',F',Q',P')=(Q^t,P^t,F,G)$ となる。
よって $C_{ij,kl}$ は $C_{kl,ji}$ に移り、二つ先では $C_{ji,lk}$ に移る。
一つ先の八次元写像は $C^*$ の入力に転置を合成したものであり、
$J$ の積を対称化したため行列対 $A,B$ は転置に関しても共変である。
従って四位置の入力側の二数 $(m,d)$ は順に
\[
 (m_V,d_V),\qquad (m_U,-d_U),\qquad (m_V,-d_V),\qquad (m_U,d_U).
\]
全二十印の合計四十個の $m_V,m_U$ が相異なり各 $d$ が非零なら八十組は相異なる。
これで異なる印への自己同型と非四倍の位置移動による同型を同時に排除できる。
四周期自体は標準Frobenius周期から従う。

印を忘れる際に二十回の代数閉体上の比較を除くこともできる。
位置を固定して二十点の有限エタール座標環を $D$ とする。$m=m_V\in D$ の二十値が
相異なるから $D=k[m]$ であり、その特性多項式 $f$ は次数20で分離的である。
$H_p$ の各表成分は $D$ 上で枠に依存せず定まり $h(m)$、$\deg h<20$ と一意に書ける。
局所的に枠を選んで求めた表は共役不変性により貼り合わさるので、分裂体を選ぶ必要はない。
従って $f$ と順序付きのこれらの係数多項式は無印の窓から有理的に求まる。
同型ならこれらは一致し、逆に一致すれば幾何学的な一根を選んだ印付きの積表が等しいので
添字保持同型を得る。このとき開始位置の移動を商にしていない。
\end{proof}

\begin{proof}[稠密開条件と全次数の合同近傍]
\path{computations/verify_twenty_mixed_bracket_invariants.py} は前節の実際の持上げから作る
四テンソルを $1009^{14}$ を法として用いる。剰余の全二十印の四位置の点を、各射影点の
一座標を固定し隣接二評価の零条件で持ち上げる。十二変数のJacobi行列に各々単元小行列があり、
持上げ後は全条件を検査する。各点を枠の第一行にし、二つの反対評価を単位に戻して
$Q,P$ を読む。持上げのスカラー部は交換子に影響しない。
各印の両線型復元の階数は両方で $18,18,36$、九語も各々基底である。
この階数消去で失う精度は最大四桁、$T$ は全て単元、$\tau$ の付値は3、
$d$ の付値は高々9、四十個の $m$ は既に $1009^3$ を法として相異なる。
これらは真の非消失であり、有限エタールな全印に条件を課しても元の十五次元成分に
空でない開集合を与える。剰余の二十点の被約性と軌道の証明は
命題\ref{prop:split-mixed-transverse-curve}の配置を用いる。

同じ三十四自由係数をこの点の値に $1009^{14}$ を法として合同に取れば、
前節の単元小行列で各々唯一の持上げを得る。各解は剰余で分裂混合窓になり、
そのWalsh符号平方行列は倍率を除き第一行 $(-2,1,1)$ の巡回行列で、余因子は $3\ne0$ である。
奇標数のWalsh判定と命題\ref{prop:unrestricted-good-reduction}より、
この近傍の\emph{各}分数体窓が幾何学的に全次数正則である。
反対階数16も前節の横断曲線から一般に成り立つ。必要ならその非消失開集合に近傍を縮められる。
自由係数を有理整数で動かせばエタール射の零次元ファイバーゆえ持上げも代数的で、
この格子は三十四次元の正規化局所成分にZariski稠密である。

最後に逆倍率も含む十九次元残余群の軌道に横断する局所切断を取る。
有理計算した積表の微分の核は正確にその軌道の接空間である。
実際、双対数環上で表が一定なら、二つの九語基底から積と単位と $C/\tau$ を保つ
一次同型が得られる。全行列環の導分は内部であり、二つの単射な積写像は
単元小行列によって双対数環でも単射なので、残りは二つの逆倍率だけである。
ここで $[C]$ の倍率を戻す一自由度も残余群に含まれる。印付き射影安定化群は自明である。
従って三十四個の自由座標のうち十九個をその軌道接空間で独立になるように固定すれば、
残る十五変数上で十五個の表成分の微分が可逆になる。
付値体の逆写像定理で十分小さい開球上ではこの十五成分が単射である。
この球内の異なる印は上の $m$ の剰余で既に区別されるから、
同じ位置で元の印を別の印に移してこれらの窓が同型になることもない。
\end{proof}
